\section{Directed edges and a guaranteed sink}\label{ch07:sec:sink}
Some connections have a direction: a one-way street, a link from one web page to another, a game won by one player against another. We record such a connection as an arrow.

\par\noindent\begin{minipage}{\linewidth}
\begin{definition}[Directed and oriented graphs]\label{ch07:def:directed}
A \emph{directed graph} consists of a set $V$ of vertices and a set of \emph{directed edges}. Each directed edge is an ordered pair $(u,v)$ of two different vertices; we write it $u\to v$ and draw it as an arrow from $u$ to $v$. An \emph{oriented graph} is a directed graph in which no two vertices are joined in both directions: $u\to v$ and $v\to u$ never both occur.
\end{definition}
\end{minipage}\par

Because the entries of an ordered pair come in a definite order (Section~\ref{ch03:sec:sets}), $u\to v$ and $v\to u$ are different directed edges. As in a graph, there are no loops $u\to u$. Some pairs of vertices may have no edge in either direction. Unlike a graph, a directed graph may have infinitely many vertices, so we say ``finite'' explicitly whenever a result needs it.

We also use $u\to w$ as a statement, saying that $(u,w)$ is a directed edge; $u\to w$ \emph{fails} when it is not. (This arrow joins two vertices; it is not the arrow of a function $f:A\to B$.) If $u\to w$, we call $w$ an \emph{out-neighbor} of $u$. The \emph{outdegree} of $u$ is the number of its out-neighbors, that is, the number of arrows leaving $u$. A \emph{sink} is a vertex with outdegree $0$: no arrow leaves it.

Walks and paths now have to follow the arrows. A \emph{directed walk} from $v_0$ to $v_k$ is a sequence of vertices $v_0,v_1,\ldots,v_k$ with $v_{i-1}\to v_i$ for every $i=1,\ldots,k$; its length is $k$. A \emph{directed path} is a directed walk with no repeated vertex. A \emph{directed cycle} is a directed walk $v_0,v_1,\ldots,v_k$ with at least one step, $k\ge1$, that returns to its start, $v_k=v_0$, and has no other repetition: the vertices $v_0,v_1,\ldots,v_{k-1}$ are all different. A directed graph is \emph{acyclic} when it has no directed cycle.

In an oriented graph, a directed cycle passes through at least three vertices, because a cycle with one step would be a loop and a cycle with two steps would need arrows in both directions. Here are two oriented graphs, each with vertices $a$, $b$, $c$:
\begin{center}
\begin{tikzpicture}[>=Stealth,every node/.style={circle,draw=ink,minimum size=6mm,inner sep=1pt},line width=.7pt]
\node(a)at(0,0){$a$};\node(b)at(1.7,0){$b$};\node(c)at(3.4,0){$c$};
\draw[->](a)--(b);\draw[->](b)--(c);
\node(a2)at(5.6,0){$a$};\node(b2)at(7.3,0){$b$};\node(c2)at(6.45,1.4){$c$};
\draw[->](a2)--(b2);\draw[->](b2)--(c2);\draw[->](c2)--(a2);
\end{tikzpicture}
\end{center}
On the left is the \emph{chain} $a\to b\to c$. It is acyclic, since every arrow leads further along the chain, and its only sink is $c$. On the right is the directed cycle $a\to b\to c\to a$. Replacing its arrow $c\to a$ by $a\to c$ gives a triangle with no directed cycle, since then every arrow goes to a later letter.

Does every directed graph have a sink? A directed graph with no vertices has none. An infinite directed graph can fail too: take the natural numbers as vertices, with an edge $n\to n+1$ for every $n$. There is no sink, yet the graph is acyclic, because the numbers strictly increase along any directed walk. Exercise~\ref{ex:7.3} asks what can happen when the graph is finite but not acyclic.

\par\noindent\begin{minipage}{\linewidth}
\begin{theorem}\label{thm:sink}
Every nonempty finite acyclic oriented graph has a sink.
\end{theorem}
\end{minipage}\par

Here \emph{nonempty} means that there is at least one vertex. The theorem promises \emph{one} sink, not that every vertex is a sink; in the chain, only $c$ is. The word ``oriented'' excludes no acyclic directed graph: a directed graph with both $u\to v$ and $v\to u$ has the directed cycle $u,v,u$. We keep the word because the conjecture in Section~\ref{ch07:sec:seymour} is about oriented graphs.

\noindent\textit{Plan.} Take a directed path that is as long as possible, and look at its last vertex. An arrow leaving that vertex would either reach a new vertex, which would give a longer path, or return to the path, which would close a directed cycle.
\begin{proof}
A directed path never repeats a vertex, so it has at most $|V|$ vertices and therefore length at most $|V|-1$. (By the pigeonhole principle, a list of more than $|V|$ vertices contains some vertex twice.) Since the graph is nonempty, it has a vertex, and a single vertex is a directed path of length $0$. So the lengths of the directed paths form a nonempty subset of $\{0,1,\ldots,|V|-1\}$, which has a largest member. Choose a directed path of this greatest length, and write it $v_0\to v_1\to\cdots\to v_k$.

Suppose, for contradiction, that its last vertex $v_k$ has an out-neighbor $w$.
\begin{itemize}
\item If $w$ is not on the path, then $v_0\to v_1\to\cdots\to v_k\to w$ is a directed path of length $k+1$, which contradicts the choice of a path of greatest length.
\item If $w$ is on the path, say $w=v_j$, then $j\ne k$, because $v_k\to v_k$ would be a loop. So $j<k$. Then $v_j\to v_{j+1}\to\cdots\to v_k\to v_j$ is a directed walk with $k-j+1$ steps that returns to its start, and its vertices $v_j,\ldots,v_k$ before the return are all different. It is a directed cycle, which contradicts acyclicity.
\end{itemize}
Both cases are impossible. Hence $v_k$ has no out-neighbor, so it is a sink.
\end{proof}

Each hypothesis has a visible job in this proof. Nonemptiness supplies a path to choose from. Finiteness guarantees a directed path of greatest length; in the directed graph on the natural numbers, every directed path can be extended. Acyclicity rules out an arrow from the last vertex back into the path.
